\begin{PSbox}[unbreakable]{obj}{Learning Objectives}

After completing this module, students will be able to:

\begin{enumerate}
\def\labelenumi{\arabic{enumi}.}
\tightlist
\item
  Extend joint distributions to $N$ random variables
\item
  Define and work with random vectors
\item
  Compute mean vectors and covariance matrices
\item
  Apply the multivariate Gaussian distribution
\item
  Perform linear transformations of random vectors
\end{enumerate}

\end{PSbox}

\PSsection{10.1}{Introduction: From 2 to $N$ Variables}

Real systems often involve many uncertain quantities simultaneously:

\begin{itemize}
\tightlist
\item
  \textbf{Sensor array:} $N$ sensors each producing a noisy measurement
\item
  \textbf{MIMO system:} $M \times N$ channel coefficients
\item
  \textbf{Parameter estimation:} $p$ unknown parameters
\item
  \textbf{Signal samples:} $N$ consecutive noise samples
\end{itemize}

We need a framework for \textbf{N-dimensional} random vectors.

\PSsection{10.2}{Joint Distributions for $N$ Random Variables}

\PSsubsection{Joint PDF}

For $X_{1},\, X_{2},\, \ldots,\, X_{n}$:\PSdisplay{f_{X_1,...,X_n}(x_1,...,x_n) \geq 0, \quad \int\cdots\int f_{X_1,...,X_n} \, dx_1 \cdots dx_n = 1}

\PSsubsection{Marginal Distributions}

Integrate (or sum) out unwanted variables:\PSdisplay{f_{X_1}(x_1) = \int\cdots\int f_{X_1,...,X_n}(x_1,...,x_n) \, dx_2 \cdots dx_n}

\PSsubsection{Conditional Distributions}

\PSdisplay{f_{X_1|X_2,...,X_n}(x_1|x_2,...,x_n) = \frac{f_{X_1,...,X_n}(x_1,...,x_n)}{f_{X_2,...,X_n}(x_2,...,x_n)}}

\PSsubsection{Mutual Independence}

$X_{1},\, \ldots,\, X_{n}$ are mutually independent if:\PSdisplay{f_{X_1,...,X_n}(x_1,...,x_n) = \prod_{i=1}^n f_{X_i}(x_i)}

\PSsection{10.3}{Random Vectors}

\begin{PSbox}[unbreakable]{def}{Definition}

A \textbf{random vector} is an ordered collection of random variables:

\PSdisplay{\mathbf{X} = \begin{bmatrix} X_1 \\ X_2 \\ \vdots \\ X_n \end{bmatrix}}

\end{PSbox}

\PSsubsection{Engineering Examples}

{\def\LTcaptype{none} % do not increment counter
\begin{longtable}[c]{@{}
  >{\centering\arraybackslash}p{(0.965\linewidth - 4\tabcolsep) * \real{0.4118}}
  >{\centering\arraybackslash}p{(0.965\linewidth - 4\tabcolsep) * \real{0.3235}}
  >{\centering\arraybackslash}p{(0.965\linewidth - 4\tabcolsep) * \real{0.2647}}@{}}
\toprule\noalign{}
\rowcolor{PSnavy}\PSth{Random Vector} & \PSth{Components} & \PSth{Context} \\
\midrule\noalign{}
\endhead
\bottomrule\noalign{}
\endlastfoot
Sensor array output & $X_{1},\,\ldots,\,X_{n}$ = sensor readings & Array processing \\
MIMO channel & $h_{1},\,\ldots,\,h_{n}$ = channel gains & Wireless communications \\
Noise vector & $N_{1},\,\ldots,\,N_{n}$ = noise samples & Signal processing \\
State vector & Position, velocity, acceleration & Kalman filtering \\
Feature vector & Extracted measurements & Pattern recognition \\
\end{longtable}
}

\PSsection{10.4}{Mean Vector}

\begin{PSbox}[unbreakable]{def}{Definition}

\PSdisplay{\boldsymbol{\mu} = E[\mathbf{X}] = \begin{bmatrix} E[X_1] \\ E[X_2] \\ \vdots \\ E[X_n] \end{bmatrix} = \begin{bmatrix} \mu_1 \\ \mu_2 \\ \vdots \\ \mu_n \end{bmatrix}}

\end{PSbox}

\PSsubsection{Properties}

\begin{itemize}
\tightlist
\item
  $E[AX + b]$ = AE[X] $+ b = A\mu + b$
\item
  Linearity extends to vectors and matrices
\end{itemize}

\PSsection{10.5}{Covariance Matrix}

\begin{PSbox}[unbreakable]{def}{Definition}

\PSdisplay{\boldsymbol{\Sigma} = E[(\mathbf{X} - \boldsymbol{\mu})(\mathbf{X} - \boldsymbol{\mu})^T]}

Element $(i,\,j)$: $\Sigma_{ij} = \operatorname{Cov}(X_{i},\, X_{j})$

\PSdisplay{\boldsymbol{\Sigma} = \begin{bmatrix} \sigma_1^2 & \text{Cov}(X_1,X_2) & \cdots & \text{Cov}(X_1,X_n) \\ \text{Cov}(X_2,X_1) & \sigma_2^2 & \cdots & \text{Cov}(X_2,X_n) \\ \vdots & & \ddots & \vdots \\ \text{Cov}(X_n,X_1) & \cdots & \cdots & \sigma_n^2 \end{bmatrix}}

\end{PSbox}

\PSsubsection{Properties}

\begin{enumerate}
\def\labelenumi{\arabic{enumi}.}
\tightlist
\item
  \textbf{Symmetric:} $\Sigma = \Sigma^{T}$
\item
  \textbf{Positive semi-definite:} $v^{T}\Sigma v \ge 0$ for all $v$
\item
  \textbf{Diagonal = variances:} $\Sigma_{ii} = \operatorname{Var}(X_{i})$
\item
  \textbf{Off-diagonal = covariances:} $\Sigma_{ij} = \operatorname{Cov}(X_{i},\, X_{j})$
\item
  \textbf{Independent components \ensuremath{\to} diagonal $\Sigma$}
\end{enumerate}

\PSsubsection{Alternative Formula}

\PSdisplay{\boldsymbol{\Sigma} = E[\mathbf{X}\mathbf{X}^T] - \boldsymbol{\mu}\boldsymbol{\mu}^T}

\PSsection{10.6}{Multivariate Gaussian Distribution}

\begin{PSbox}[unbreakable]{def}{Definition}

$X \sim N(\mu,\, \Sigma)$ has PDF:

\PSdisplay{f_\mathbf{X}(\mathbf{x}) = \frac{1}{(2\pi)^{n/2}|\boldsymbol{\Sigma}|^{1/2}} \exp\left(-\frac{1}{2}(\mathbf{x}-\boldsymbol{\mu})^T \boldsymbol{\Sigma}^{-1}(\mathbf{x}-\boldsymbol{\mu})\right)}

where $\lvert \Sigma\rvert$ = determinant of $\Sigma$.

\end{PSbox}

\PSsubsection{Key Properties}

\begin{enumerate}
\def\labelenumi{\arabic{enumi}.}
\tightlist
\item
  \textbf{Marginals are Gaussian:} Any subset of components is jointly Gaussian
\item
  \textbf{Conditionals are Gaussian:} $X_{1} \mid X_{2} = x_{2}$ is Gaussian
\item
  \textbf{Linear transformations:} $Y$ = AX $+ b \sim N(A\mu +b,\, A\Sigma A^{T})$
\item
  \textbf{Uncorrelated = Independent} (unique to Gaussian!)
\item
  \textbf{Fully specified} by $\mu$ and $\Sigma$ (only 2 parameters needed per variable plus correlations)
\end{enumerate}

\PSsubsection{Why So Important?}

\begin{itemize}
\tightlist
\item
  Central Limit Theorem: sums of many RVs \ensuremath{\to} multivariate Gaussian
\item
  Thermal noise vectors are multivariate Gaussian
\item
  Analytically tractable (closed-form marginals, conditionals)
\item
  Optimal filters (Kalman, Wiener) assume Gaussian
\end{itemize}

\PSsubsection{$2D$ Case (Bivariate Gaussian)}

For $n = 2$ with $\mu = [0,\,0]^{T},\, \sigma_{1} = \sigma_{2} = 1$:

\PSdisplay{f(x,y) = \frac{1}{2\pi\sqrt{1-\rho^2}} \exp\left(-\frac{x^2 - 2\rho xy + y^2}{2(1-\rho^2)}\right)}

Contours are \textbf{ellipses} (tilted when $\rho \ne 0$).

\PSsection{10.7}{Linear Transformations of Random Vectors}

\PSsubsection{Transformation $Y$ = AX $+ b$}

If $X$ is a random vector with mean $\mu_{X}$ and covariance $\Sigma_{X}$:

\PSdisplay{E[\mathbf{Y}] = A\boldsymbol{\mu}_X + \mathbf{b}}\PSdisplay{\boldsymbol{\Sigma}_Y = A\boldsymbol{\Sigma}_X A^T}

\PSsubsection{Engineering Applications}

\textbf{Beamforming:} Output $y = w^{T}X$ (weight vector applied to sensor array)

\begin{itemize}
\tightlist
\item
  $E[y] = w^{T}\mu$
\item
  $\operatorname{Var}(y) = w^{T}\Sigma w$
\end{itemize}

\textbf{Whitening:} Find matrix $W$ such that $W\Sigma W^{T}$ = I (decorrelates data)

\begin{itemize}
\tightlist
\item
  $W = \Sigma^{-1/2}$
\end{itemize}

\textbf{Principal Component Analysis:} Rotate to diagonalize $\Sigma$

\PSsection{10.8}{MATLAB Examples}

\begin{PSbox}[unbreakable]{ex}{Example 1: Generating Multivariate Gaussian}

\tcbset{PScodemode/.style={unbreakable}}

\begin{Shaded}
\begin{Highlighting}[]
\CommentTok{\%\% Generate and Visualize 2D Gaussian Random Vector}
\VariableTok{mu} \OperatorTok{=}\NormalTok{ [}\FloatTok{1}\OperatorTok{;} \FloatTok{2}\NormalTok{]}\OperatorTok{;}
\VariableTok{Sigma} \OperatorTok{=}\NormalTok{ [}\FloatTok{4}\OperatorTok{,} \FloatTok{2}\OperatorTok{;} \FloatTok{2}\OperatorTok{,} \FloatTok{3}\NormalTok{]}\OperatorTok{;}  \CommentTok{\% Covariance matrix}

\VariableTok{N} \OperatorTok{=} \FloatTok{5000}\OperatorTok{;}
\VariableTok{X} \OperatorTok{=} \VariableTok{mvnrnd}\NormalTok{(}\VariableTok{mu}\OperatorTok{\textquotesingle{},} \VariableTok{Sigma}\OperatorTok{,} \VariableTok{N}\NormalTok{)}\OperatorTok{;}

\VariableTok{figure}\OperatorTok{;}
\VariableTok{scatter}\NormalTok{(}\VariableTok{X}\NormalTok{(}\OperatorTok{:,}\FloatTok{1}\NormalTok{)}\OperatorTok{,} \VariableTok{X}\NormalTok{(}\OperatorTok{:,}\FloatTok{2}\NormalTok{)}\OperatorTok{,} \FloatTok{5}\OperatorTok{,} \SpecialStringTok{\textquotesingle{}b\textquotesingle{}}\OperatorTok{,} \SpecialStringTok{\textquotesingle{}.\textquotesingle{}}\NormalTok{)}\OperatorTok{;}
\VariableTok{hold} \VariableTok{on}\OperatorTok{;}
\VariableTok{plot}\NormalTok{(}\VariableTok{mu}\NormalTok{(}\FloatTok{1}\NormalTok{)}\OperatorTok{,} \VariableTok{mu}\NormalTok{(}\FloatTok{2}\NormalTok{)}\OperatorTok{,} \SpecialStringTok{\textquotesingle{}r+\textquotesingle{}}\OperatorTok{,} \SpecialStringTok{\textquotesingle{}MarkerSize\textquotesingle{}}\OperatorTok{,} \FloatTok{15}\OperatorTok{,} \SpecialStringTok{\textquotesingle{}LineWidth\textquotesingle{}}\OperatorTok{,} \FloatTok{3}\NormalTok{)}\OperatorTok{;}
\VariableTok{xlabel}\NormalTok{(}\SpecialStringTok{\textquotesingle{}X\_1\textquotesingle{}}\NormalTok{)}\OperatorTok{;} \VariableTok{ylabel}\NormalTok{(}\SpecialStringTok{\textquotesingle{}X\_2\textquotesingle{}}\NormalTok{)}\OperatorTok{;}
\VariableTok{title}\NormalTok{(}\SpecialStringTok{\textquotesingle{}Bivariate Gaussian Samples\textquotesingle{}}\NormalTok{)}\OperatorTok{;}
\VariableTok{axis} \VariableTok{equal}\OperatorTok{;} \VariableTok{grid} \VariableTok{on}\OperatorTok{;}

\CommentTok{\% Verify statistics}
\VariableTok{fprintf}\NormalTok{(}\SpecialStringTok{\textquotesingle{}Mean: [\%.2f, \%.2f] (theory: [\%.1f, \%.1f])\textbackslash{}n\textquotesingle{}}\OperatorTok{,} \VariableTok{mean}\NormalTok{(}\VariableTok{X}\NormalTok{)}\OperatorTok{,} \VariableTok{mu}\OperatorTok{\textquotesingle{}}\NormalTok{)}\OperatorTok{;}
\VariableTok{fprintf}\NormalTok{(}\SpecialStringTok{\textquotesingle{}Cov matrix:\textbackslash{}n\textquotesingle{}}\NormalTok{)}\OperatorTok{;} \VariableTok{disp}\NormalTok{(}\VariableTok{cov}\NormalTok{(}\VariableTok{X}\NormalTok{))}\OperatorTok{;}
\end{Highlighting}
\end{Shaded}

\end{PSbox}

\begin{PSbox}[unbreakable]{ex}{Example 2: Linear Transformation}

\tcbset{PScodemode/.style={unbreakable}}

\begin{Shaded}
\begin{Highlighting}[]
\CommentTok{\%\% Transform: Y = AX + b}
\VariableTok{mu\_x} \OperatorTok{=}\NormalTok{ [}\FloatTok{0}\OperatorTok{;} \FloatTok{0}\NormalTok{]}\OperatorTok{;} \VariableTok{Sigma\_x} \OperatorTok{=} \VariableTok{eye}\NormalTok{(}\FloatTok{2}\NormalTok{)}\OperatorTok{;}   \CommentTok{\% Standard normal}
\VariableTok{A} \OperatorTok{=}\NormalTok{ [}\FloatTok{2}\OperatorTok{,} \FloatTok{1}\OperatorTok{;} \OperatorTok{{-}}\FloatTok{1}\OperatorTok{,} \FloatTok{3}\NormalTok{]}\OperatorTok{;}  \VariableTok{b} \OperatorTok{=}\NormalTok{ [}\FloatTok{1}\OperatorTok{;} \OperatorTok{{-}}\FloatTok{2}\NormalTok{]}\OperatorTok{;}

\VariableTok{N} \OperatorTok{=} \FloatTok{10000}\OperatorTok{;}
\VariableTok{X} \OperatorTok{=} \VariableTok{mvnrnd}\NormalTok{(}\VariableTok{mu\_x}\OperatorTok{\textquotesingle{},} \VariableTok{Sigma\_x}\OperatorTok{,} \VariableTok{N}\NormalTok{)}\OperatorTok{\textquotesingle{};}
\VariableTok{Y} \OperatorTok{=} \VariableTok{A}\OperatorTok{*}\VariableTok{X} \OperatorTok{+} \VariableTok{b}\OperatorTok{;}

\CommentTok{\% Theoretical}
\VariableTok{mu\_y\_theory} \OperatorTok{=} \VariableTok{A}\OperatorTok{*}\VariableTok{mu\_x} \OperatorTok{+} \VariableTok{b}\OperatorTok{;}
\VariableTok{Sigma\_y\_theory} \OperatorTok{=} \VariableTok{A}\OperatorTok{*}\VariableTok{Sigma\_x}\OperatorTok{*}\VariableTok{A}\OperatorTok{\textquotesingle{};}

\VariableTok{fprintf}\NormalTok{(}\SpecialStringTok{\textquotesingle{}E[Y] theory: [\%.1f; \%.1f], sim: [\%.2f; \%.2f]\textbackslash{}n\textquotesingle{}}\OperatorTok{,} \OperatorTok{...}
    \VariableTok{mu\_y\_theory}\OperatorTok{,} \VariableTok{mean}\NormalTok{(}\VariableTok{Y}\OperatorTok{,}\FloatTok{2}\NormalTok{))}\OperatorTok{;}
\VariableTok{fprintf}\NormalTok{(}\SpecialStringTok{\textquotesingle{}Σ\_Y theory:\textbackslash{}n\textquotesingle{}}\NormalTok{)}\OperatorTok{;} \VariableTok{disp}\NormalTok{(}\VariableTok{Sigma\_y\_theory}\NormalTok{)}\OperatorTok{;}
\VariableTok{fprintf}\NormalTok{(}\SpecialStringTok{\textquotesingle{}Σ\_Y simulated:\textbackslash{}n\textquotesingle{}}\NormalTok{)}\OperatorTok{;} \VariableTok{disp}\NormalTok{(}\VariableTok{cov}\NormalTok{(}\VariableTok{Y}\OperatorTok{\textquotesingle{}}\NormalTok{))}\OperatorTok{;}
\end{Highlighting}
\end{Shaded}

\end{PSbox}

\begin{PSbox}[breakable]{ex}{Example 3: Sensor Array with Correlated Noise}

\tcbset{PScodemode/.style={unbreakable}}

\begin{Shaded}
\begin{Highlighting}[]
\CommentTok{\%\% 4{-}Element Sensor Array with Spatially Correlated Noise}
\VariableTok{n\_sensors} \OperatorTok{=} \FloatTok{4}\OperatorTok{;}
\VariableTok{sigma2} \OperatorTok{=} \FloatTok{1}\OperatorTok{;}      \CommentTok{\% Noise power per sensor}
\VariableTok{rho} \OperatorTok{=} \FloatTok{0.6}\OperatorTok{;}       \CommentTok{\% Adjacent sensor correlation}

\CommentTok{\% Build covariance matrix (exponential decay model)}
\VariableTok{Sigma} \OperatorTok{=} \VariableTok{zeros}\NormalTok{(}\VariableTok{n\_sensors}\NormalTok{)}\OperatorTok{;}
\KeywordTok{for} \VariableTok{i} \OperatorTok{=} \FloatTok{1}\OperatorTok{:}\VariableTok{n\_sensors}
    \KeywordTok{for} \VariableTok{j} \OperatorTok{=} \FloatTok{1}\OperatorTok{:}\VariableTok{n\_sensors}
        \VariableTok{Sigma}\NormalTok{(}\VariableTok{i}\OperatorTok{,}\VariableTok{j}\NormalTok{) }\OperatorTok{=} \VariableTok{sigma2} \OperatorTok{*} \VariableTok{rho}\OperatorTok{\^{}}\VariableTok{abs}\NormalTok{(}\VariableTok{i}\OperatorTok{{-}}\VariableTok{j}\NormalTok{)}\OperatorTok{;}
    \KeywordTok{end}
\KeywordTok{end}

\VariableTok{fprintf}\NormalTok{(}\SpecialStringTok{\textquotesingle{}Noise covariance matrix:\textbackslash{}n\textquotesingle{}}\NormalTok{)}\OperatorTok{;} \VariableTok{disp}\NormalTok{(}\VariableTok{Sigma}\NormalTok{)}\OperatorTok{;}

\CommentTok{\% Generate noise samples}
\VariableTok{N\_samp} \OperatorTok{=} \FloatTok{10000}\OperatorTok{;}
\VariableTok{L} \OperatorTok{=} \VariableTok{chol}\NormalTok{(}\VariableTok{Sigma}\OperatorTok{,} \SpecialStringTok{\textquotesingle{}lower\textquotesingle{}}\NormalTok{)}\OperatorTok{;}
\VariableTok{noise} \OperatorTok{=} \VariableTok{L} \OperatorTok{*} \VariableTok{randn}\NormalTok{(}\VariableTok{n\_sensors}\OperatorTok{,} \VariableTok{N\_samp}\NormalTok{)}\OperatorTok{;}

\CommentTok{\% Signal present in direction: s = [1; 1; 1; 1] (broadside)}
\VariableTok{signal} \OperatorTok{=} \FloatTok{3}\OperatorTok{;}
\VariableTok{s} \OperatorTok{=} \VariableTok{ones}\NormalTok{(}\VariableTok{n\_sensors}\OperatorTok{,} \FloatTok{1}\NormalTok{) }\OperatorTok{/} \VariableTok{sqrt}\NormalTok{(}\VariableTok{n\_sensors}\NormalTok{)}\OperatorTok{;}
\VariableTok{received} \OperatorTok{=} \VariableTok{signal} \OperatorTok{*} \VariableTok{s} \OperatorTok{+} \VariableTok{noise}\OperatorTok{;}  \CommentTok{\% N\_sensors x N\_samp}

\CommentTok{\% Beamformer output: y = w\textquotesingle{}*received (w = s for matched filter)}
\VariableTok{w} \OperatorTok{=} \VariableTok{s}\OperatorTok{;}
\VariableTok{y} \OperatorTok{=} \VariableTok{w}\OperatorTok{\textquotesingle{}} \OperatorTok{*} \VariableTok{received}\OperatorTok{;}
\VariableTok{SNR\_out} \OperatorTok{=} \VariableTok{signal}\OperatorTok{\^{}}\FloatTok{2} \OperatorTok{/}\NormalTok{ (}\VariableTok{w}\OperatorTok{\textquotesingle{}*}\VariableTok{Sigma}\OperatorTok{*}\VariableTok{w}\NormalTok{)}\OperatorTok{;}
\VariableTok{fprintf}\NormalTok{(}\SpecialStringTok{\textquotesingle{}Output SNR = \%.2f (\%.1f dB)\textbackslash{}n\textquotesingle{}}\OperatorTok{,} \VariableTok{SNR\_out}\OperatorTok{,} \FloatTok{10}\OperatorTok{*}\VariableTok{log10}\NormalTok{(}\VariableTok{SNR\_out}\NormalTok{))}\OperatorTok{;}
\end{Highlighting}
\end{Shaded}

\end{PSbox}

\PSsection{10.9}{Practice Problems}

\begin{PSbox}[unbreakable]{prob}{Problem 1}

Random vector $X = [X_{1},\, X_{2},\, X_{3}]^{T}$ with $\mu = [1,\, 0,\, -1]^{T}$, and covariance matrix $\Sigma = \begin{bmatrix} 4 & 1 & 0 \\ 1 & 2 & -1 \\ 0 & -1 & 3 \end{bmatrix}$.

\begin{enumerate}
\def\labelenumi{(\alph{enumi})}
\tightlist
\item
  Find $\operatorname{Var}(X_{1} + X_{2} + X_{3})$. (b) Find $\operatorname{Cov}(X_{1},\, X_{2}+X_{3})$. (c) Are $X_{1}$ and $X_{3}$ uncorrelated?
\end{enumerate}

\PSsolution{} 

\begin{enumerate}
\def\labelenumi{(\alph{enumi})}
\tightlist
\item
  $\operatorname{Var}(X_{1}+X_{2}+X_{3}) = 1^{T}\Sigma 1 = 4+2+3 + 2(1) + 2(0) + 2(-1) = 9 + 0 = 9$\PSbr{}Wait: $= \sum_{ij}$ all entries $= 4+1+0+1+2+(-1)+0+(-1)+3 = 9$
\item
  $\operatorname{Cov}(X_{1},\, X_{2}+X_{3}) = \operatorname{Cov}(X_{1},\,X_{2}) + \operatorname{Cov}(X_{1},\,X_{3}) = 1 + 0 = 1$
\item
  $\operatorname{Cov}(X_{1},\,X_{3}) = 0$, so yes, $X_{1}$ and $X_{3}$ are uncorrelated.
\end{enumerate}

\end{PSbox}

\begin{PSbox}[unbreakable]{prob}{Problem 2}

$X \sim N(0,\, \Sigma)$ with $\Sigma = \begin{bmatrix} 1 & \rho \\ \rho & 1 \end{bmatrix}$. Find the conditional distribution of $X_{1} \mid X_{2} = x_{2}$.

\PSsolution{} For bivariate Gaussian:\PSbr{}$\displaystyle X_{1} \mid X_{2} = x_{2} \sim N\left(\mu_{1} + \rho \left(\frac{\sigma_{1}}{\sigma_{2}}\right)(x_{2}-\mu_{2}),\, \sigma_{1}^{2}(1-\rho^{2})\right)$\PSbr{}$= N(\rho x_{2},\, 1-\rho^{2})$

The conditional mean is a linear function of $x_{2}$, and conditional variance is reduced by factor $(1-\rho^{2})$.

\end{PSbox}

\begin{PSbox}[unbreakable]{prob}{Problem 3}

Write MATLAB code to generate a $3D$ Gaussian vector with given $\mu$ and $\Sigma$, compute the sample covariance, and verify.

\PSsolution{} 

\tcbset{PScodemode/.style={unbreakable}}

\begin{Shaded}
\begin{Highlighting}[]
\VariableTok{mu} \OperatorTok{=}\NormalTok{ [}\FloatTok{1}\OperatorTok{;} \FloatTok{0}\OperatorTok{;} \OperatorTok{{-}}\FloatTok{1}\NormalTok{]}\OperatorTok{;}
\VariableTok{Sigma} \OperatorTok{=}\NormalTok{ [}\FloatTok{4} \FloatTok{1} \FloatTok{0}\OperatorTok{;} \FloatTok{1} \FloatTok{2} \OperatorTok{{-}}\FloatTok{1}\OperatorTok{;} \FloatTok{0} \OperatorTok{{-}}\FloatTok{1} \FloatTok{3}\NormalTok{]}\OperatorTok{;}
\VariableTok{N} \OperatorTok{=} \FloatTok{50000}\OperatorTok{;}
\VariableTok{X} \OperatorTok{=} \VariableTok{mvnrnd}\NormalTok{(}\VariableTok{mu}\OperatorTok{\textquotesingle{},} \VariableTok{Sigma}\OperatorTok{,} \VariableTok{N}\NormalTok{)}\OperatorTok{;}
\VariableTok{fprintf}\NormalTok{(}\SpecialStringTok{\textquotesingle{}Sample mean: \textquotesingle{}}\NormalTok{)}\OperatorTok{;} \VariableTok{disp}\NormalTok{(}\VariableTok{mean}\NormalTok{(}\VariableTok{X}\NormalTok{)}\OperatorTok{\textquotesingle{}}\NormalTok{)}\OperatorTok{;}
\VariableTok{fprintf}\NormalTok{(}\SpecialStringTok{\textquotesingle{}Sample cov:\textbackslash{}n\textquotesingle{}}\NormalTok{)}\OperatorTok{;} \VariableTok{disp}\NormalTok{(}\VariableTok{cov}\NormalTok{(}\VariableTok{X}\NormalTok{))}\OperatorTok{;}
\end{Highlighting}
\end{Shaded}

\emph{Next Module: {Module 11 — Central Limit Theorem}}

\end{PSbox}
